\FloatBarrier
\subsection{Axes IR}
\label{sec:axes-ir}
Axes IR is TEPL's intermediate representation of tensor computations. It
describes how operations use tensor axes and compute output values.
Translating operation definitions into this shared representation lets the
compiler track axis relationships across operations, infer replacement
attributes, and check whether completed rewrites preserve the result.

\FloatBarrier
\subsubsection{Symbolic Axis Groups}
\label{sec:axis-groups}
A symbolic axis group represents an ordered sequence of axes without fixing
its length. One group can, for example, describe all contraction axes of an
operation. This representation borrows arbitrary-rank grouping from
TensorRight's aggregated axes~\cite{tensorright}. TensorRight represents an
aggregated axis as a set of named axes; TEPL uses an ordered sequence and
tracks position mappings from groups to tensor axes for attribute construction.

\begin{figuretext}{figures/Symbolic axis groups.pdf}{Ordered, arbitrary-length axis groups preserve operand pairings and determine output order.}{fig:axis-groups}{Four arbitrary-length groups for batch, left-free, contraction, and right-free
axes are connected to operand position lists and an output-axis construction.}
TEPL declares a group with \code{bind}, associating it with axis positions in
one or more operands. When two operands bind the same group, their selected
axes are paired in list order. Attributes determine those positions, allowing
the same definition to describe different ranks and layouts. The two selections
must have equal length and equal sizes at corresponding positions. Matching
sizes alone do not establish pairing: the binding specifies which coordinates
correspond.

Figure~\ref{fig:axis-groups} illustrates arbitrary-length groups, one concrete
dot binding, and output construction. In that instance, the contraction axes
appear in different physical orders in the two operands. Selecting \code{[2,3]}
in \code{X} and \code{[5,2]} in \code{Y} nevertheless identifies the same
ordered contraction group. Reducing it leaves the ordered output
$B\catseq I\catseq J$.


We use four related notions. A tensor's ordered axes form an interface
$P=\langle a_0,\ldots,a_{n-1}\rangle$ in dimension order, together with established
mappings from groups into $P$. Axis labels denote logical coordinates, rather
than sizes. A group $G=\langle g_0,\ldots,g_{m-1}\rangle$ is an ordered
sequence whose order can differ from dimension order. A position sequence
$p=\langle p_0,\ldots,p_{m-1}\rangle$, with $0\leq p_j<n$, locates that
group by $g_j=a_{p_j}$. Positions are zero-based. Finally, $\gath(P,p)$
selects axes in the order specified by $p$, so a mapping records
\begin{equation}
 \gath(P,p)=G.
 \label{eq:mapping}
\end{equation}
\end{figuretext}

The following \code{dot_general} definition declares four attributes and their
axis bindings. Dtype and value semantics are omitted from this excerpt:
\clearpage
\begin{codeblock}[H]
\begin{lstlisting}[language=TEPL]
op dot_general(lhs: tensor, rhs: tensor) -> tensor {
  attrs {
    lhs_batching_dimensions: index[];
    rhs_batching_dimensions: index[];
    lhs_contracting_dimensions: index[];
    rhs_contracting_dimensions: index[];
  }
  axes {
    let lb = attrs.lhs_batching_dimensions;
    let rb = attrs.rhs_batching_dimensions;
    let lc = attrs.lhs_contracting_dimensions;
    let rc = attrs.rhs_contracting_dimensions;
    let lf = ordered_complement(all_axes(lhs), lb ++ lc);
    let rf = ordered_complement(all_axes(rhs), rb ++ rc);
    bind B { lhs: lb, rhs: rb };
    bind K { lhs: lc, rhs: rc };
    bind I { lhs: lf };
    bind J { rhs: rf };
    yield B ++ I ++ J;
  }
}
\end{lstlisting}
\end{codeblock}
Here, \code{all_axes(t)} lists all positions of tensor \code{t}, and \code{++}
concatenates sequences. The builtin \code{ordered_complement} selects remaining
positions in operand order. The axis \code{yield} specifies output order.

As a second concrete instance, smaller than the one in
Figure~\ref{fig:axis-groups}, consider:
\begin{codeblock}[H]
\begin{lstlisting}
lhs:    [i, k0, b, k1]
rhs:    [k1, j, k0, b]
output: [b, i, j]

lhs_batching_dimensions = [2]
rhs_batching_dimensions = [3]
lhs_contracting_dimensions = [1, 3]
rhs_contracting_dimensions = [2, 0]
\end{lstlisting}
\end{codeblock}
The batch binding pairs position $2$ of \code{lhs} with position $3$ of
\code{rhs}. The contraction lists pair $(1,2)$ for $k_0$ and $(3,0)$ for
$k_1$. Remaining axes give $I=\langle i\rangle$ and $J=\langle j\rangle$;
$K$ is reduced and therefore omitted from the output. In mapping notation,
\begin{equation}
\begin{aligned}
P_{\mathrm{lhs}}&=\langle i,k_0,b,k_1\rangle,
 &p_{\mathrm{lhs}}&=\langle1,3\rangle,\\
P_{\mathrm{rhs}}&=\langle k_1,j,k_0,b\rangle,
 &p_{\mathrm{rhs}}&=\langle2,0\rangle,\\
G_{\mathrm{lhs}}&=G_{\mathrm{rhs}}=K=\langle k_0,k_1\rangle.
\end{aligned}
\end{equation}
Consequently,
\begin{equation}
 \gath(P_{\mathrm{lhs}},p_{\mathrm{lhs}})
 =K=\gath(P_{\mathrm{rhs}},p_{\mathrm{rhs}}).
\end{equation}
Although $K$ has two axes here, neither group length nor position values need
be known at rule-compilation time. Grouping preserves the axes separately and
in order; it does not flatten them into one dimension. Groups can be empty
where an operation permits it, such as an unbatched dot's batch group. Their
identities are scoped so that unrelated coordinates with equal extents remain
distinct. Mappings let the compiler recover where each logical group occurs
and use those positions as replacement attributes.

\FloatBarrier
\subsubsection{Expressing Tensor Computations}
\label{sec:value-expressions}
The \code{axes} block defines relationships and output order, while
\code{semantics} defines each output element's value. Axes IR distinguishes
\emph{value expressions}, which compute scalar tensor elements, from
\emph{index expressions}, which calculate the indices used to read them.
Ordinary operand references use coordinates established by bindings. Explicit
index expressions handle calculated accesses, including slices and
convolutions. \code{Read} connects the two kinds of expression by reading an
operand at explicitly calculated indices. A \code{semantics} block returns a
value expression with \code{yield}.

For example, elementwise addition without broadcasting binds every axis of
\code{lhs} to the corresponding axis of \code{rhs}, preserving that group:
\begin{codeblock}[H]
\begin{lstlisting}[language=TEPL]
axes {
  bind S {
    lhs: all_axes(lhs),
    rhs: all_axes(rhs)
  };
  yield S;
}
semantics {
  yield Add(lhs, rhs);
}
\end{lstlisting}
\end{codeblock}
The shared binding requires equal rank and matching sizes. At each coordinate
in $S$, the operand references read corresponding elements and \code{Add}
produces their sum. Table~\ref{tab:value-expressions} lists the value primitives.
In this language, \code{Reduce} means summation; it does not imply support
for other reductions such as maximum.

\begin{table}[htbp]
\caption{Value-expression primitives.}
\label{tab:value-expressions}
\begin{tabularx}{\linewidth}{@{}lY@{}}\toprule
Expression & Meaning\\\midrule
\code{input} & Read an operand element using its axis bindings.\\
\code{Const(c)} & Produce a signed integer constant.\\
\code{Add(a, b)} & Add scalar values.\\
\code{Product(a, b)} & Multiply scalar values.\\
\code{Reduce(G, e)} & Sum $e$ over all coordinate combinations of group $G$.\\
\code{Read(input, indices)} & Read an operand at explicitly calculated indices.\\\bottomrule
\end{tabularx}
\end{table}

Index expressions produce an index sequence with one entry for each accessed
axis. Table~\ref{tab:index-expressions} gives their forms. A slice, for example,
uses
\begin{codeblock}[H]
\begin{lstlisting}[language=TEPL]
Read(input,
  IndexOffset(attrs.starts,
    IndexScale(attrs.strides, Coords(O))))
\end{lstlisting}
\end{codeblock}
Here $O$ is the output group. Per axis, the index expression computes
\code{starts + strides * output_indices}. The \code{Read} expression retrieves
the selected value. The index calculation determines the location, and the
enclosing value expression determines the scalar result.

\begin{table}[htbp]
\caption{Index-expression primitives.}
\label{tab:index-expressions}
\begin{tabularx}{\linewidth}{@{}lY@{}}\toprule
Expression & Meaning\\\midrule
\code{Coords(G)} & Current coordinates of $G$ in group order.\\
\code{IndexAdd(a, b)} & Add index sequences elementwise.\\
\code{Concat(a, b)} & Concatenate index sequences.\\
\code{IndexScale(attrs.strides, indices)} & Scale each index by its corresponding entry.\\
\code{IndexOffset(attrs.starts, indices)} & Add the corresponding offset to each index.\\\bottomrule
\end{tabularx}
\end{table}

Value and index expressions share the groups declared in \code{axes}.
Bindings determine participating operand elements; references bound to the
same group use corresponding coordinates. A reduction ranges over its group
while holding the remaining output coordinates fixed. Because a group can
have variable length, one expression can sum over different numbers of
dimensions. The value semantics completing the previous dot definition are:
\begin{codeblock}[H]
\begin{lstlisting}[language=TEPL]
semantics {
  yield Reduce(K, Product(lhs, rhs));
}
\end{lstlisting}
\end{codeblock}
The product reads coordinates determined by $B,I,J,K$, and the reduction sums
over $K$. Output coordinates belong to $B,I,J$, in the declared output order.
The smaller concrete example therefore computes
\begin{equation}
 \operatorname{out}[b,i,j]
 =\sum_{k_0,k_1}\operatorname{lhs}[i,k_0,b,k_1]
                    \operatorname{rhs}[k_1,j,k_0,b].
\end{equation}

These primitives describe elementwise operations, transpose, summation,
and dot products. Slices and supported convolutions calculate input indices;
broadcasting adds output axes while reusing input values along them. Output
shape remains separate from element computation: a tensor filled with zeros,
for example, still needs a specified shape.

\FloatBarrier
\subsubsection{Reasoning About Composed Operations}
\label{sec:composition}
Value semantics provide a common language for tensor rewrites. Different
operations lower to scalar addition, multiplication, and summation, so laws
over these primitives can relate expressions built from different tensor
operations. TEPL normalizes both sides under the arithmetic laws permitted
by the selected policy.

Consider distributing a generalized dot over elementwise addition without
broadcasting, with identical dot descriptors on both sides:
\begin{equation}
 \operatorname{dot\_general}(X+Y,Z)
 \;\longleftrightarrow\;
 \operatorname{dot\_general}(X,Z)+\operatorname{dot\_general}(Y,Z).
\end{equation}
Addition's local value expression is \code{Add(X, Y)} under group $S$.
The dot's local expression is \code{Reduce(K, Product(lhs, Z))} under its own
batch, contraction, and free groups. Substituting the arithmetic expression
for \code{lhs} gives the expected shape of the computation, but does not
establish how positions in $S$ correspond to $K$ or retain their pairing with
\code{Z}'s contraction axes.

Symbolic axis mapping supplies this missing correspondence. The dot's
contraction attribute selects positions in addition's output. Since addition
preserves axis order, those positions identify corresponding axes in both
\code{X} and \code{Y}. The compiler records this relationship, and similarly
connects batch and free axes, even when group lengths are unknown. With fixed
output coordinates and indexed reads abbreviated, the two expressions become
\begin{equation}
 \sum_K(X+Y)Z
 \qquad\text{and}\qquad
 \sum_K XZ+\sum_K YZ.
\end{equation}
In the value language they are
\begin{codeblock}[H]
\begin{lstlisting}[language=TEPL]
Reduce(K, Product(Add(X, Y), Z))

Add(Reduce(K, Product(X, Z)),
    Reduce(K, Product(Y, Z)))
\end{lstlisting}
\end{codeblock}
Here, operand names abbreviate reads under aligned bindings, and $K$ is the
corresponding contraction group on both sides. Distribution of multiplication
and summation over addition brings the expressions to a common form when
the arithmetic policy permits those laws. Section~\ref{sec:axis-mapping}
explains how the required mappings are represented, composed, and refined.
